\documentclass{article}
\usepackage{amsmath,amssymb}
\begin{document}
\section{Complete product proof for the original theta entire function}
This appendix proves the genus-zero factorization used in the
counterfactual workbench. It is a proof of the classical special case,
included to make the application's analytic dependencies explicit.

\subsection{A sublinear growth bound and its zeros}
Let $F$ be entire, $F(0)\ne0$, and suppose that for some
$0<\alpha<1$ and constants $C,R_0>0$,
\[
\log\max_{|z|\le R}|F(z)|\le CR^\alpha\qquad(R\ge R_0).
\tag{GF1}
\]
The workbench's original $F$ satisfies this with $\alpha=3/4$:
its explicitly proved bound $\log M_F(R)=O(\sqrt R\log R)$
implies (GF1) because $\log R/R^{1/4}$ is bounded for $R\ge1$.
Let the nonzero zeros be $a_j$, repeated with multiplicity, and let
$n(r)$ count those with $|a_j|\le r$.

We first prove the required counting formula. Choose a circle
$|z|=R$ containing no zero. There are finitely many zeros in its
closed disk, since otherwise an accumulation point would force the
nonzero entire function to vanish identically. Divide $F$ by their
linear factors; the resulting holomorphic nonvanishing function
on a neighborhood of the disk has a holomorphic logarithm there.
Indeed its logarithmic derivative is holomorphic on the disk and
has a primitive, as follows by termwise integration of its Taylor
series on concentric smaller disks. The exponential of that primitive
has the same logarithmic derivative; their quotient is constant,
giving the logarithm after choosing the constant at zero.
The real part of this logarithm has its value at zero as its circle
mean, by integrating its Taylor series.
For $|a|<R$,
\[
\frac1{2\pi}\int_0^{2\pi}\log|Re^{it}-a|\,dt=\log R.
\tag{GF2}
\]
To see this, expand
$\log(1-ae^{-it}/R)=-\sum_{\ell\ge1}(a/R)^\ell e^{-i\ell t}/\ell$,
which converges uniformly on the circle; each term has zero mean.
Taking the logarithm of the removed factors now gives
\[
\frac1{2\pi}\int_0^{2\pi}\log|F(Re^{it})|\,dt-\log|F(0)|
=\sum_{|a_j|<R}\log\frac{R}{|a_j|}.
\tag{GF3}
\]
For any sufficiently large $r$, choose $R\in(2r,3r)$ containing
no zero. All terms are nonnegative and every zero with
$|a_j|\le r$ contributes at least $\log2$. Equations (GF1)--(GF3)
therefore imply $n(r)\le C_1r^\alpha$, for a fixed positive
$C_1$ and all sufficiently large $r$.

Choose a positive $r_0$ smaller than the modulus of every zero;
such a radius exists since $F(0)\ne0$. After enlarging $C_1$,
the inequality $n(r)\le C_1r^\alpha$ holds for every $r\ge r_0$.
The sum over successive dyadic annuli gives
\[
\sum_j\frac1{|a_j|}
\le\sum_{\ell\ge0}
\frac{n(2^{\ell+1}r_0)}{2^\ell r_0}
\le C_1\,2^\alpha r_0^{\alpha-1}
\sum_{\ell\ge0}2^{(\alpha-1)\ell}<\infty.
\tag{GF4}
\]
The inequalities can overcount zeros and still give the stated
upper bound. If the zero set is empty, every sum is zero.

\subsection{The product and a lower bound on selected circles}
Define the literal product
\[
P(z)=\prod_j(1-z/a_j).
\tag{GF5}
\]
On a fixed disk $|z|\le A$, all sufficiently late factors have
$|z/a_j|\le1/2$. The power series logarithm satisfies
$|\log(1-z/a_j)|\le2|z/a_j|$, so (GF4) makes its tail uniformly
convergent there. The finite preceding factors and that exponential
tail define an entire function, locally uniformly independent of
the ordering. Its zeros and their multiplicities are exactly the
$a_j$, and $P(0)=1$. Thus
\[
H(z)=F(z)/P(z)
\tag{GF6}
\]
extends through every zero to a nonvanishing entire function.
An entire primitive of $H'/H$, obtained by integrating its entire
power series, gives an entire logarithm $h$ with $H=e^h$ after
adjusting its constant.

We prove $h$ constant without assuming any growth bound for
$1/P$ away from the following selected circles. Let $t$ be
sufficiently large, put $N_t=n(4t)$, and set
$\delta_t=t/[4(N_t+1)]$. Delete from $[t,2t]$ each open interval
of radius $\delta_t$ about a number $|a_j|\le4t$.
Their total length is at most $2N_t\delta_t<t/2$.
Hence a radius $R_t\in[t,2t]$ remains with
$|R_t-|a_j||\ge\delta_t$ for every such zero. For $|z|=R_t$,
\[
|1-z/a_j|\ge\frac{\delta_t}{4t}
=\frac1{16(N_t+1)}\qquad(|a_j|\le4t).
\tag{GF7}
\]
For the remaining zeros, $|z/a_j|\le1/2$, and the same logarithm
estimate used for (GF5) gives
$\log|1-z/a_j|\ge-2|z|/|a_j|$.
Another dyadic estimate, retaining only $|a_j|>4t$, gives
\[
\sum_{|a_j|>4t}\frac1{|a_j|}
\le\sum_{\ell\ge0}
\frac{n(2^{\ell+1}4t)}{2^\ell4t}
\le C_2t^{\alpha-1},
\tag{GF8}
\]
where $C_2=C_1\,2^\alpha4^{\alpha-1}/(1-2^{\alpha-1})$
is positive and fixed. Summing the logarithms of (GF7) and the
uniformly convergent tail proves, on the entire selected circle,
\[
\log|P(z)|
\ge-N_t\log(16(N_t+1))-2R_tC_2t^{\alpha-1}
\ge-C_3t^\alpha\log t
\tag{GF9}
\]
for another fixed constant $C_3>0$ and all large $t$.
This estimate keeps every zero, including repetitions.

\subsection{The zero-free factor is constant}
Combining (GF1), (GF6), and (GF9) shows that
$\Re h(z)\le C_4t^\alpha\log t$ on $|z|=R_t$.
Write $h(z)=\sum_{\ell\ge0}d_\ell z^\ell$ and put
$B_t=C_4t^\alpha\log t$, enlarged if necessary so that the bound
holds on every selected circle under consideration.
The Taylor series, uniformly convergent on a circle, gives
\[
d_\ell R_t^\ell=
\frac1{\pi}\int_0^{2\pi}\Re h(R_te^{i\theta})
e^{-i\ell\theta}\,d\theta\qquad(\ell\ge1).
\tag{GF10}
\]
The integral of the constant $B_t$ against
$e^{-i\ell\theta}$ is zero. Because
$B_t-\Re h\ge0$ on the circle, its mean and the mean-value
identity yield
\[
|d_\ell|R_t^\ell
\le\frac1\pi\int_0^{2\pi}(B_t-\Re h(R_te^{i\theta}))\,d\theta
=2(B_t-\Re h(0)).
\tag{GF11}
\]
Since $R_t\ge t$ and $\alpha<1$, the right side divided by
$R_t^\ell$ tends to zero for every integer $\ell\ge1$.
Thus every such $d_\ell$ is zero. It follows that $H$ is
constant. Evaluating at zero now proves the full factorization
\[
F(z)=F(0)\prod_j(1-z/a_j).
\tag{GF12}
\]
This also covers finitely many zeros and the empty zero set.

\subsection{Logarithmic coefficients with their full multiplicities}
Group equal zeros in (GF12), writing distinct reciprocals
$\beta_j=a_j^{-1}$ and multiplicities $m_j$. Equation (GF4)
becomes $\sum_jm_j|\beta_j|<\infty$.
Their supremum is finite. Choose a sufficiently small positive
$r$ such that $r|\beta_j|\le1/2$ for all $j$.
The derivative series
$\sum_jm_j\beta_j/(1-\beta_j z)$ converges uniformly on
$|z|\le r$, by domination with $2\sum_jm_j|\beta_j|$.
It is the negative logarithmic derivative of (GF12), by
termwise differentiation of its locally uniformly convergent
logarithm, or by integrating the uniformly convergent derivative
series and matching its value at zero. Expanding each denominator
and interchanging absolutely convergent sums gives
\[
-\frac{F'(z)}{F(z)}
=\sum_j\frac{m_j\beta_j}{1-\beta_jz}
=\sum_{n\ge0}
\left(\sum_jm_j\beta_j^{n+1}\right)z^n .
\tag{GF13}
\]
Every series interchange is justified by
$\sum_jm_j|\beta_j|/(1-r|\beta_j|)<\infty$.
Equations (GF12)--(GF13) prove exactly the product and power-sum
input used in the workbench and (TC16), with the unchanged
$F(0)=g(1/2)$.

\end{document}
